% !TeX root = ../../Book2.tex
% 纲要标题：### 14.4 相对与带边界横截性
% 纲要位置：几何学计划纲要.md，第 1563 行。
% 正文状态：初稿；尚未正式审核。

% 纲要内容开始
% 
% * 局部扰动和有限覆盖
% * 保持指定区域不变
% * 边界限制映射的横截条件
% * 同伦族与配边的生成
% 
% ---
% 
% 纲要内容结束

\section{相对与带边界横截性}
\label{g2:sec:1404}

参数定理要求先有一个总体横截的映射族。对任意给定映射，怎样造出这样的族？在目标坐标中加一个小向量，就能补足缺少的方向；在源端乘截断函数，又能把改变限制在指定区域。将这个局部操作逐块实施，便得到保持已给数据的一般位置结果。边界问题还需检查沿边界本身是否有足够的方向，最后才能把一族逆像拼成配边。

\subsection{局部扰动和有限覆盖}
\label{g2:subsec:1404:01}

\begin{lemma}[坐标内的横截扰动]\label{g2:lem:local-transverse-perturbation}
	设 \(M,N\) 是无边界光滑流形，\(S\subset N\) 是嵌入子流形，\(f:M\to N\) 光滑。设 \(K\subset U\subset M\)，其中 \(K\) 紧，\(U\) 开，\(f(U)\) 落在一张目标坐标图内。则可在 \(U\) 内选择紧支扰动，使所得光滑映射 \(g\) 在 \(K\) 上横截于 \(S\)，并在 \(M\setminus U\) 上等于 \(f\)。扰动可以在任意预先指定的有限阶 \(C^r\) 意义下任意小。
\end{lemma}
\begin{proof}
	记目标坐标为 \(y:V\to\mathbb R^n\)。第一卷第12章第3节的光滑截断结论保证存在 \(\rho:M\to[0,1]\)，在 \(K\) 的某邻域上为 \(1\)，且支集紧含于 \(U\)。取充分小的参数球 \(B\subset\mathbb R^n\)，定义
	\[
		F(a,x)=\begin{cases}
			y^{-1}\bigl(y(f(x))+\rho(x)a\bigr),&x\in U,\\
			f(x),&x\notin\operatorname{supp}\rho.
		\end{cases}
	\]
	两个表达在交叠处一致；\(f(\operatorname{supp}\rho)\) 紧含于 \(V\)，故可统一缩小 \(B\)，使平移后的点仍在坐标像内。在 \(\rho>0\) 的区域，关于参数的微分在坐标中为 \(\rho(x)I_n\)，所以总体映射 \(F\) 是淹没。\cref{g2:thm:parametric-transversality}应用于 \(B\times\{\rho>0\}\)，说明几乎每个任意小的 \(a\) 都使 \(F_a\) 在这里横截于 \(S\)，尤其在 \(K\) 上横截。

	令 \(g=F_a\)。在固定紧支集上，\(F_a\) 的各阶坐标导数连续依赖于 \(a\)，并在 \(a=0\) 时等于 \(f\) 的相应导数。因此让 \(a\) 足够小，即可满足任意给定的有限阶误差界。对非紧源上的强拓扑，紧支集只与局部有限控制族中的有限多个紧块相交，同一结论仍适用。
\end{proof}

这个构造说明参数为什么应在目标坐标中取值：源端原本可能缺少的法向方向，由 \(a\) 直接补上。截断函数的过渡区不必横截；下一步会用覆盖把所有尚未处理的点纳入某个 \(\rho=1\) 的区域。

\ProseParagraphBreak

先设 \(M\) 紧，\(S\) 闭。取有限个紧块 \(K_i\subset U_i\)，其内部覆盖 \(M\)，每个 \(f(U_i)\) 落在一张目标图中。依次应用引理，使第 \(i\) 步在 \(K_i\) 上横截。已经处理的 \(K_1\cup\cdots\cup K_{i-1}\) 是紧集；\cref{g2:prop:transversality-open-on-compact}保证，只要本步足够小，就保留那里的横截性。再把总误差分配成有限份，最后得到任意接近 \(f\) 的处处横截映射。这里保持的是先前取得的满秩条件，先前的函数值可以继续改变。

\subsection{保持指定区域不变}
\label{g2:subsec:1404:02}

若映射在某个区域已经用于拼接，那里常常必须逐点不变。例如同伦的两个端点已经给定，不能在“取得一般位置”时顺便更换端点。相对横截性把这种要求直接列入结论。

\begin{theorem}[相对横截逼近]\label{g2:thm:relative-transversality}
	设 \(M,N\) 是有限维、Hausdorff、第二可数的无边界光滑流形，\(S\subset N\) 是闭嵌入子流形，\(A\subset M\) 闭。若光滑映射 \(f:M\to N\) 在 \(A\) 的一个开邻域上横截于 \(S\)，则对每个有限整数 \(r\ge1\) 及 \(f\) 的每个强 \(C^r\) 邻域 \(\mathcal U\)，存在光滑映射 \(g\in\mathcal U\)，使 \(g\pitchfork S\)，并使 \(g=f\) 在 \(A\) 的某个开邻域上成立。\(g\) 还可选为与 \(f\) 光滑同伦，且同伦在这个较小邻域上不变。\(A=\varnothing\) 时即得到普通横截逼近。
\end{theorem}
\begin{proof}
	设 \(f\) 在开集 \(U_0\supset A\) 上横截。由第一卷第4章第3—4节的正规性与覆盖收缩，选开集 \(W,V\)，使
	\[
		A\subset W\subset\overline W\subset V\subset\overline V\subset U_0.
	\]
	可先缩小 \(\mathcal U\)，使其中每个映射在 \(\overline V\) 上横截。具体做法是在 \(\overline V\) 附近取局部有限的紧坐标块，在每块上用紧集横截开放性给出正的 \(C^1\) 容差；强拓扑允许同时要求这些逐块不等式。因而接下来只需处理 \(M\setminus V\)。

	对这个闭集取可数、在整个 \(M\) 上局部有限的相对紧开集族 \(U_i\)，每个 \(\overline{U_i}\) 避开 \(\overline W\)，且 \(f(\overline{U_i})\) 含于一个目标坐标图。取紧块 \(K_i\subset U_i\)，使其内部覆盖 \(M\setminus V\)。这里使用的是第一卷第12章第3节构造单位分解时的局部有限收缩；尤其局部有限性在 \(W\) 的边缘也成立。用逐块 \(C^0\) 容差保证后续映射仍留在所选目标图内。

	从 \(f_0=f\) 开始，在第 \(i\) 步应用\cref{g2:lem:local-transverse-perturbation}，得到 \(f_i\)，使其在 \(K_i\) 上横截，改动支集紧含于 \(U_i\)。同时让改动足够小，以保留 \(K_1,\ldots,K_{i-1}\) 上的横截性，并满足原邻域的误差要求。下面说明无穷多步中如何保留这些要求，避免只在每个有限阶段正确。

	在固定的源、目标坐标块中用各阶导数的上确界测量差。若原邻域在第 \(\ell\) 块给出容差 \(\varepsilon_\ell\)，令第 \(i\) 步在它碰到的各块上耗用不到 \(2^{-i-1}\varepsilon_\ell\)。第 \(i\) 步完成后，紧集开放性给出 \(\delta_i>0\)：以后只要在覆盖 \(K_i\) 的有限坐标块上与 \(f_i\) 的 \(C^1\) 差小于 \(\delta_i\)，便仍横截。对每个 \(j>i\)，额外要求第 \(j\) 次改动在这些块上的大小小于 \(2^{-(j-i+1)}\delta_i\)。总和小于 \(\delta_i\)。一次改动有紧支，所遇原控制块有限，先前的 \(K_i\) 也只有有限多个；所以每一步只是对参数添加有限多个严格上界，可以同时满足。

	每一点的某邻域只遇到有限多个 \(U_i\)，因此 \(f_i\) 在该邻域最终不再改变，定义出光滑映射 \(g\)。上述误差和保证 \(g\in\mathcal U\) 并在每个 \(K_i\) 上横截；在 \(\overline V\) 上的横截性由最初缩小的邻域保证。这些集合覆盖 \(M\)，故 \(g\pitchfork S\)。所有改动支集避开 \(\overline W\)，于是 \(g=f\) 于 \(W\)。

	最后，在各次坐标平移中把选定参数 \(a_i\) 换成 \(t a_i\)，按相同次序组合，\(0\le t\le1\)。各步可一并选得对所有 \(t\) 满足前述坐标内的误差预算，所以组合有定义。源端的局部有限性保证每点附近只有有限次组合，从而给出联合光滑的 \(H(t,x)\)，满足 \(H(0,x)=f(x)\)、\(H(1,x)=g(x)\)，且在 \(W\) 上恒等于 \(f\)。这个同伦本身的每张切片不必横截。
\end{proof}

\begin{corollary}[余维较大的集合可以避开]\label{g2:cor:avoid-high-codimension}
	设 \(M,N\) 是无边界光滑流形，\(S\subset N\) 是闭嵌入子流形，且 \(\dim M<\operatorname{codim}_N S\)。每个光滑映射 \(M\to N\) 都可任意光滑逼近为像避开 \(S\) 的映射，并保持同伦类；若原映射在闭集 \(A\) 的某邻域上避开 \(S\)，可在较小邻域内保持它不变。
\end{corollary}
\begin{proof}
	相对横截逼近给出 \(g\pitchfork S\)。若 \(g(x)\in S\)，则 \(\mathrm{T}_xM\to\mathrm{T}_{g(x)}N/\mathrm{T}_{g(x)}S\) 必须满射，与维数不等式矛盾。因此 \(g^{-1}(S)=\varnothing\)。任意 \(C^\infty\) 邻域只需满足构成它的一组有限阶基本条件，选足够大的有限 \(r\) 即可。
\end{proof}

例如三维空间中的一条闭路可以作任意小的扰动避开给定点；对一张曲面也成立，因为点的余维为3。若障碍物是一张曲面，余维只有1，上述维数论证便不再给出避开结论。相交数将在后续章节解释何时交点反而不能全部消除。

\subsection{边界限制映射的横截条件}
\label{g2:subsec:1404:03}

\begin{theorem}[带边界的横截逆像]\label{g2:thm:boundary-transverse-preimage}
	设 \(M^m\) 是光滑带边界流形，\(N\) 是无边界光滑流形，\(S\subset N\) 是无边界、余维为 \(c\) 的嵌入子流形。若光滑映射 \(f:M\to N\) 以及限制 \(f|_{\partial M}\) 都横截于 \(S\)，则 \(Z=f^{-1}(S)\) 若非空，就是 \(m-c\) 维带边界嵌入子流形，且
	\[
		\partial Z=Z\cap\partial M,\qquad
		\mathrm{T}_xZ=(\mathrm{d}f_x)^{-1}(\mathrm{T}_{f(x)}S).
	\]
	在 \(x\in\partial Z\) 处还有
	\[
		\mathrm{T}_x(\partial Z)=\mathrm{T}_xZ\cap\mathrm{T}_x(\partial M),\qquad
		\mathrm{T}_xZ+\mathrm{T}_x(\partial M)=\mathrm{T}_xM.
	\]
\end{theorem}
\begin{proof}
	内部点由\cref{g2:thm:transverse-preimage}处理。边界点处取 \(S\) 的局部定义函数 \(q:V\to\mathbb R^c\)，即 \(S\cap V=q^{-1}(0)\)，且 \(\ker\mathrm{d}q=\mathrm{T}S\) 沿 \(S\) 成立。令 \(h=q\circ f\)。边界横截条件等价于 \(\mathrm{d}(h|_{\partial M})\) 满射，这正是第一卷第11章第4节“带边界的正则水平集”所需的条件。

	为看清边界为何保留下来，在半空间坐标 \((x^1,\ldots,x^m)\)、\(x^m\ge0\) 中，满射性使前 \(m-1\) 个方向中的某 \(c\) 列给出可逆子式。重排这些坐标后，以
	\[
		\bigl(h^1(x),\ldots,h^c(x),x^{c+1},\ldots,x^{m-1},x^m\bigr)
	\]
	作新坐标。逆函数定理适用于坐标函数的局部光滑延拓，且最后一个坐标仍为 \(x^m\)，故局部逆保持半空间。\(Z\) 在这些坐标中就是前 \(c\) 项为零的半空间切片。切片边界、切空间及最后一个向内方向分别给出陈述中的各个等式。包含映射在这些图中为坐标嵌入，图册与子空间拓扑相容。
\end{proof}

\begin{example}\label{g2:exm:boundary-transversality-failure}
	在上半平面 \(M=\{(x,y):y\ge0\}\) 上令 \(f(x,y)=y+x^2\)，取 \(S=\{0\}\subset\mathbb R\)。由于 \(\partial_y f=1\)，\(f\) 处处为淹没；然而 \(f^{-1}(0)=\{(0,0)\}\)，并不是预期的一维带边界流形。缺少的条件是 \(f|_{\partial M}(x)=x^2\) 在零点处的横截性。改成 \(g(x,y)=x\)，逆像为向上的半直线，才具有一维内部及一个边界点。完整切空间含向内法向，而边界限制只能使用沿边界的方向，两项检查因此有不同作用。
\end{example}

相对逼近也有相同的边界版本：在\cref{g2:thm:relative-transversality}中允许源 \(M\) 带边界，要求 \(f\) 与 \(f|_{\partial M}\) 在固定区域都横截，则可使最终 \(g\) 与 \(g|_{\partial M}\) 同时横截。证明中的坐标扰动在半空间上仍有意义；在 \(\rho>0\) 处，总体映射限制到 \(B\times\operatorname{Int}M\) 和 \(B\times\partial M\) 都因参数方向而成为淹没。对这两张无边界源的映射分别应用参数定理，避开两个零测异常集的并，即可同时满足条件。边界上的满射还蕴含完整切空间上的满射。紧块开放性分别应用于内部的完整导数和边界切向导数，此后的有限覆盖或局部有限构造完全沿用上述误差选择。

\subsection{同伦族与配边的生成}
\label{g2:subsec:1404:04}

一条同伦使每个点的像随时间移动；若整条同伦横截于 \(S\)，所有相交点的运动轨迹就组成高一维的流形。它的边界记录起点和终点的逆像。这是由一般位置构造几何不变量的基本方式。

\begin{definition}[配边]\label{g2:def:unoriented-cobordism}
	设 \(Z_0,Z_1\) 是紧无边界 \(k\) 维光滑流形。若存在紧 \((k+1)\) 维光滑带边界流形 \(W\)，以及微分同胚 \(\partial W\cong Z_0\amalg Z_1\)，就称 \(Z_0,Z_1\)\term[peibian]{配边}[cobordant]。此处暂不指定定向；有向版本还需规定边界的诱导定向。
\end{definition}

\begin{theorem}[同伦的横截逆像]\label{g2:thm:transverse-homotopy-cobordism}
	设 \(M^m\) 是紧无边界光滑流形，\(N\) 是无边界光滑流形，\(S\subset N\) 是闭、余维 \(c\le m\) 的嵌入子流形。若两个光滑映射 \(f_0,f_1:M\to N\) 均横截于 \(S\)，且连续同伦，则存在光滑同伦 \(H:[0,1]\times M\to N\)，在两个端点的时间邻域分别恒等于 \(f_0,f_1\)，并且 \(H\pitchfork S\)。于是
	\[
		W=H^{-1}(S),\qquad
		\partial W=\bigl(\{0\}\times f_0^{-1}(S)\bigr)
		\amalg\bigl(\{1\}\times f_1^{-1}(S)\bigr)
	\]
	给出两个逆像之间的配边。
\end{theorem}
\begin{proof}
	从一条连续同伦 \(F\) 出发，取光滑非降函数 \(\lambda:[0,1]\to[0,1]\)，在0附近等于0、在1附近等于1。把 \(F\) 换成 \(F(\lambda(t),x)\)，所得同伦在两端时间邻域中等于光滑的 \(f_0,f_1\)。\cref{g2:thm:relative-smooth-approximation}适用于这个带边界源，给出保持较小端点邻域的光滑同伦 \(F'\)。

	在这些端点邻域内，\(F'\) 以及它的边界限制都横截于 \(S\)，因为 \(f_0,f_1\) 已横截。应用刚证明的带边界相对横截逼近，固定再小一些的端点邻域，得到所需的 \(H\)。\cref{g2:thm:boundary-transverse-preimage}给出 \(W\) 的流形结构及边界公式。最后，\(S\) 闭，所以 \(W\) 是紧空间 \([0,1]\times M\) 的闭子集，因而紧；这一点正是配边定义中的紧性要求。
\end{proof}

\begin{example}\label{g2:exm:zeros-appear-in-pairs}
	在 \(S^1\) 上考虑 \(f_a(e^{\mathrm i\theta})=\cos\theta-a\)。当 \(|a|<1\) 时，零点有两个且都是正则零点；当 \(|a|>1\) 时没有零点。沿 \(a\) 从0变到2的总体函数，其参数导数等于 \(-1\)，故总体逆像是光滑一维流形，虽然 \(a=1\) 的单张切片有临界零点。两个零点的轨迹在该临界时刻接合并消失。这说明总体横截足以构造配边，却不要求每个时刻的切片都横截。
\end{example}

一般地，当 \(m=c\) 时，横截逆像是有限点集。上述配边是紧一维流形，其边界点数为偶数。可以直接用有限图解释这件事：取有限条闭坐标弧，其相对内部覆盖 \(W\)，遇到边界点时使用半区间图。把所有弧端点和 \(\partial W\) 放入有限集 \(P\)。任一坐标弧删去 \(P\) 后的各开段，在 \(W\setminus P\) 中既开又闭，故其连通分支就是有限条这样的开段。这把 \(W\) 写成有限图：内部顶点的度数为2，边界顶点的度数为1。由各边恰有两个端点，全部顶点度数之和为边数的两倍，故度数为1的顶点有偶数个。因此 \(f_0^{-1}(S)\) 与 \(f_1^{-1}(S)\) 的点数模2相同。后续相交数理论将在加入定向后，把这个奇偶计数细化为带符号计数。
